% ================= 第 2 章: GL2(Qp) 的序表示 =================

\section{$\GL_2(\Qp)$ 的序表示}

在本部分, 在关于抛物线诱导的预备知识之后, 我们定义引言中的酉 $\GL_2(\Qp)$-Banach 空间 $B(\sigma_p)$.
固定 $L$ 为 $\Qp$ 中 $\Qp$ 的有限扩张.

\subsection{\mbox{}}\label{sec:2.1}

我们首先讨论关于抛物线诱导的一些一般性考虑.
记 $\widehat{T}$ 为 $\Qp$ 上参数化 $T(\Qp)$ 的局部解析 (或连续, 二者等价) 特征标 $\chi_1 \otimes \chi_2$ 的刚性解析空间.
此空间同构于乘积 $(\mathcal{W} \times \mathbb{G}_m)^{2}$, 其中 $\mathcal{W}$ 是 $\Qp$ 上参数化 $\Zp^{\times}$ 的局部解析特征标的刚性解析空间 (\cite[\S4.4]{ref:24}).
因此 $\widehat{T}(L)$ 的一个元素就是一个局部解析特征标 $\chi_1 \otimes \chi_2 : T(\Qp) \to L^{\times}$.

\begin{definition}\label{def:2.1.1}
称 $\widehat{T}(L)$ 的特征标 $\chi_1 \otimes \chi_2$ 为经典权 $k$ 的, 其中 $k \ge 2$ 为整数, 若特征标 $\chi_2/\chi_1 : \Qp^{\times} \to L^{\times}$ 是一个局部常值特征标与代数特征标 $z \mapsto z^{k-2}$ 的乘积.
称 $\widehat{T}(L)$ 的特征标 $\chi_1 \otimes \chi_2$ 为经典权的, 若存在整数 $k \ge 2$ 使得 $\chi_1 \otimes \chi_2$ 是经典权 $k$ 的.
\end{definition}

若 $d$ 为正整数或 $0$, 记 $\mathcal{C}^{\mathrm{lp},\le d}(\Zp,L)$ (resp.\ $\mathcal{C}^{\mathrm{an}}(\Zp,L)$, resp.\ $\mathcal{C}^{0}(\Zp,L)$) 为局部多项式次数 $\le d$ (resp.\ 局部解析, resp.\ 连续) 的函数 $f : \Zp \to L$ 的 $L$-向量空间, 配备最细的局部凸拓扑 (resp.\ \cite{ref:44} 中定义的紧型拓扑, resp.\ 与范数 $\mathrm{Sup}\,|f(z)|$ 相关联的拓扑).
当 $d = 0$ 时, 也记 $\mathcal{C}^{\mathrm{lc}}(\Zp,L) := \mathcal{C}^{\mathrm{lp},\le 0}(\Zp,L)$ 为 $\Zp$ 上局部常值函数的空间.

若 $\chi_1 \otimes \chi_2 : T(\Qp) \to L^{\times}$ 是经典权 $k$ (resp.\ 局部解析, resp.\ 连续) 的特征标, 记
\[
\pind{\mathrm{lp},\le k-2}_{B(\Qp)}^{\GL_2(\Qp)} \chi_1 \otimes \chi_2
\qquad
\left(\text{resp.}\ \anind_{B(\Qp)}^{\GL_2(\Qp)} \chi_1 \otimes \chi_2,\ \text{resp.}\ \cind_{B(\Qp)}^{\GL_2(\Qp)} \chi_1 \otimes \chi_2\right)
\]
为局部多项式次数 $\le k-2$ (resp.\ \cite[\S5,6]{ref:44} 意义下的局部解析, resp.\ \cite{ref:43} 意义下的连续) 抛物线诱导, 即函数 $F : \GL_2(\Qp) \to L$ 的 $L$-向量空间 $V$, 其中 $g \mapsto \det(g)^{-1}F(g)$ 局部多项式次数 $\le k-2$ (resp.\ 局部解析, resp.\ 连续), 且满足:
\[
F\!\left(\begin{pmatrix} a & b \\ 0 & d \end{pmatrix}g\right) = \chi_1(a)\chi_2(d)F(g),
\]
配备由 $(g \cdot F)(g') = F(g'g)$ 给出的 $\GL_2(\Qp)$ 左作用.
此空间自然配备一个局部凸拓扑 (在局部代数情形为最细拓扑, 在局部解析情形为紧型拓扑 (精确定义参见 \cite{ref:44}), 在连续情形为与范数 $\mathrm{Sup}_{g \in \GL_2(\Zp)}|f(g)|$ 相关联的拓扑), 使得映射 $\GL_2(\Qp) \times V \to V$, $(g,F) \mapsto g \cdot F$ 连续.

映射:
\[
F \in V \mapsto \left(z \mapsto F\!\begin{pmatrix} 0 & 1 \\ -1 & z \end{pmatrix}\right)
\]
将 $V$ 等同于局部多项式次数 $\le k-2$ (resp.\ 局部解析, resp.\ 连续) 的函数 $f : \Qp \to L$ 的 $L$-向量空间, 其中 $(\chi_2\chi_1^{-1})(z)f(1/z)$ 在 $z = 0$ 的邻域可延拓为次数 $\le k-2$ 的多项式 (resp.\ 解析, resp.\ 连续) 函数.
此时 $\GL_2(\Qp)$ 的作用写作:
\begin{equation}
\label{eq:1}
\begin{pmatrix} a & b \\ c & d \end{pmatrix} \cdot f(z) = \chi_1(ad-bc)(\chi_2\chi_1^{-1})(-cz+a)\,f\!\left(\frac{dz-b}{-cz+a}\right),
\end{equation}
而映射:
\[
f \mapsto \left(z \mapsto f(pz),\ z \mapsto (\chi_2\chi_1^{-1})(z)f(1/z)\right)
\]
给出与 $\mathcal{C}^{\mathrm{lp},\le k-2}(\Zp,L)^{2}$ (resp.\ $\mathcal{C}^{\mathrm{an}}(\Zp,L)^{2}$, resp.\ $\mathcal{C}^{0}(\Zp,L)^{2}$) 的拓扑同构.
若 $k = 2$, 也记 $\lcind_{B(\Qp)}^{\GL_2(\Qp)} \chi_1 \otimes \chi_2 := \pind{\mathrm{lp},\le 0}_{B(\Qp)}^{\GL_2(\Qp)} \chi_1 \otimes \chi_2$.
它就是通常的 $1 \otimes \chi_2\chi_1^{-1}$ 被 $\chi_1$ 扭曲后的光滑诱导.

我们需要考虑前述局部解析诱导中不被 $\GL_2(\Qp)$ 保持、但被 $\mathfrak{gl}_2(\Qp)$ 和 $B(\Qp)$ 保持的子空间.
现在引入它们.
对 $\chi_1 \otimes \chi_2 \in \widehat{T}(L)$ 和整数 $d \ge 0$, 记
\[
\bigl(\pind{}_{B(\Qp)}^{\GL_2(\Qp)} \chi_1 \otimes \chi_2\bigr)(\Qp)^{\mathrm{lp},\le d}
\]
为 $\anind_{B(\Qp)}^{\GL_2(\Qp)} \chi_1 \otimes \chi_2$ 中支撑在 $\Qp$ 内紧、且局部多项式次数 $\le d$ 的上述函数 $f$ 的子空间.
若 $d = 0$, 也记 $\bigl(\pind{}_{B(\Qp)}^{\GL_2(\Qp)} \chi_1 \otimes \chi_2\bigr)(\Qp)^{\mathrm{lc}}$.
设:
\[
\bigl(\pind{}_{B(\Qp)}^{\GL_2(\Qp)} \chi_1 \otimes \chi_2\bigr)(\Qp)^{\mathrm{lp}} := \bigcup_{d \ge 0} \bigl(\pind{}_{B(\Qp)}^{\GL_2(\Qp)} \chi_1 \otimes \chi_2\bigr)(\Qp)^{\mathrm{lp},\le d}.
\]
给这些空间配备最细的局部凸拓扑.
每个空间 $\bigl(\pind{}_{B(\Qp)}^{\GL_2(\Qp)} \chi_1 \otimes \chi_2\bigr)(\Qp)^{\mathrm{lp},\le d}$ 都是 $B(\Qp)$-不变的, 且它们的并还是 $\mathfrak{gl}_2(\Qp)$-不变的.
包含:
\[
\bigl(\pind{}_{B(\Qp)}^{\GL_2(\Qp)} \chi_1 \otimes \chi_2\bigr)(\Qp)^{\mathrm{lc}} \subset \bigl(\pind{}_{B(\Qp)}^{\GL_2(\Qp)} \chi_1 \otimes \chi_2\bigr)(\Qp)^{\mathrm{lp}}
\]
因此诱导一个 $(\mathfrak{gl}_2(\Qp),B(\Qp))$-等变态射:
\begin{equation}
\label{eq:2}
U_{\mathfrak{gl}_2(\Qp)} \otimes_{U_{\mathfrak{b}(\Qp)}} \bigl(\pind{}_{B(\Qp)}^{\GL_2(\Qp)} \chi_1 \otimes \chi_2\bigr)(\Qp)^{\mathrm{lc}} \to \bigl(\pind{}_{B(\Qp)}^{\GL_2(\Qp)} \chi_1 \otimes \chi_2\bigr)(\Qp)^{\mathrm{lp}},
\end{equation}
其中 $U_{\mathfrak{gl}_2(\Qp)}$ (resp.\ $U_{\mathfrak{b}(\Qp)}$) 是 $\mathfrak{gl}_2(\Qp)$ (resp.\ $\mathfrak{b}(\Qp)$) 的包络代数.

下面的两个引理是经典的.
为方便读者, 我们给出它们的证明.

\begin{lemma}\label{lem:2.1.2}
若 $\chi_1 \otimes \chi_2$ 不是经典权的, 则 \eqref{eq:2} 是同构.
\end{lemma}

\begin{proof}
设 $w_i \in L$ 为特征标 $\chi_i$ ($i = 1,2$) 的权, 记 $w := w_2 - w_1$.
要求 $\chi_1 \otimes \chi_2$ 不是经典权的, 等价于要求 $w$ 不是正整数或 $0$.
若 $X = \begin{pmatrix} \alpha & \beta \\ \gamma & \delta \end{pmatrix} \in \mathfrak{gl}_2(\Qp)$ 且 $f \in \bigl(\pind{}_{B(\Qp)}^{\GL_2(\Qp)} 1\bigr)(\Qp)^{\mathrm{lp}}$ (如前作为 $\Qp$ 上的函数), 对 \eqref{eq:1} 求导得到关于 $X \cdot f$ 的如下公式:
\begin{equation}
\label{eq:3}
(X \cdot f)(z) = \bigl(w_1(\alpha+\delta) + w(\alpha-\gamma z)\bigr)f(z) + \bigl(-\beta + (\delta-\alpha)z + \gamma z^{2}\bigr)f'(z),
\end{equation}
其中 $f'(z)$ 是 $f$ 的导数.
设 $L[z]$ 为系数在 $L$ 中的 $z$ 的多项式的 $L$-向量空间.
公式 \eqref{eq:3} 定义了 $\mathfrak{gl}_2(\Qp)$ 在 $L[z]$ 上的一个表示, 它是 $\mathfrak{gl}_2(\Qp)$ 的 Verma 反步模的一个例子.
常值函数子空间 $L \subset L[z]$ 被 $\mathfrak{n}(\Qp)$ 零化, 而 $\mathfrak{t}(\Qp)$ 通过特征标 $\chi_2|\,| \otimes \chi_1|\,|^{-1}$ 的导数作用在 $L$ (看作 $1$ 维 $L$-向量空间) 上, 即 $\begin{pmatrix} \alpha & 0 \\ 0 & \delta \end{pmatrix} \mapsto w_2\alpha + w_1\delta$.
$\mathfrak{gl}_2(\Qp)$ 在 $L[z]$ 上的作用诱导一个 $U_{\mathfrak{gl}_2(\Qp)}$-模态射:
\begin{equation}
\label{eq:4}
U_{\mathfrak{gl}_2(\Qp)} \otimes_{U_{\mathfrak{b}(\Qp)}} L \to L[z],
\end{equation}
而这是 Verma 模理论中的一个经典结果 (可直接用 \eqref{eq:3} 验证): 在 $w$ 不是正整数或 $0$ 的假设下, \eqref{eq:4} 是同构.
若 $\Omega$ 是 $\Qp$ 的紧开集, 记 $L[z]|_{\Omega}$ 为 $\Qp \to L$ 的函数空间, 它们在 $\Omega$ 上是多项式且在其补集上为零.
由 \eqref{eq:3} 可见 $L[z]|_{\Omega}$ 是 $\bigl(\pind{}_{B(\Qp)}^{\GL_2(\Qp)} \chi_1 \otimes \chi_2\bigr)(\Qp)^{\mathrm{lp}}$ 中被 $\mathfrak{gl}_2(\Qp)$ 稳定的子空间.
若 $f \in L[z]$, 设 $f|_{\Omega}$ 为 $L[z]|_{\Omega}$ 中在 $\Omega$ 上等于多项式 $f$、在 $\Qp \setminus \Omega$ 上为 $0$ 的元素.
映射 $f \mapsto f|_{\Omega}$ 诱导 $\mathfrak{gl}_2(\Qp)$-模同构 $L[z] \xrightarrow{\sim} L[z]|_{\Omega}$.
若 $L|_{\Omega}$ 表示 $L[z]|_{\Omega}$ 中在 $\Omega$ 上为常值的函数构成的子空间, 则有交换图:
\[
\begin{CD}
U_{\mathfrak{gl}_2(\Qp)} \otimes_{U_{\mathfrak{b}(\Qp)}} L @>>> L[z]\\
@V{\mathrm{id} \otimes (f \mapsto f|_{\Omega})}VV @VV{f \mapsto f|_{\Omega}}V\\
U_{\mathfrak{gl}_2(\Qp)} \otimes_{U_{\mathfrak{b}(\Qp)}} L|_{\Omega} @>>> L[z]|_{\Omega}
\end{CD}
\]
其中上面的水平箭头和两个竖直箭头都是同构.
因此下面的水平箭头也是同构.
另一方面, 有 $\mathfrak{gl}_2(\Qp)$-模 (resp.\ $\mathfrak{b}(\Qp)$-模) 同构:
\[
\varinjlim_{\{\Omega_i\}} \bigoplus_i L[z]|_{\Omega_i} \xrightarrow{\sim} \bigl(\pind{}_{B(\Qp)}^{\GL_2(\Qp)} \chi_1 \otimes \chi_2\bigr)(\Qp)^{\mathrm{lp}}
\]
\[
\left(\text{resp.}\ \varinjlim_{\{\Omega_i\}} \bigoplus_i L|_{\Omega_i} \xrightarrow{\sim} \bigl(\pind{}_{B(\Qp)}^{\GL_2(\Qp)} \chi_1 \otimes \chi_2\bigr)(\Qp)^{\mathrm{lc}}\right)
\]
其中归纳极限取遍 $\Qp$ 的互不相交紧开集的所有有限族 $\{\Omega_i\}$.
但由上述结果, 自然箭头:
\[
U_{\mathfrak{gl}_2(\Qp)} \otimes_{U_{\mathfrak{b}(\Qp)}} \varinjlim_{\{\Omega_i\}} \bigoplus_i L|_{\Omega_i} \to \varinjlim_{\{\Omega_i\}} \bigoplus_i L[z]|_{\Omega_i}
\]
是同构, 由此推出引理.
\end{proof}

\begin{lemma}\label{lem:2.1.3}
若 $\chi_1 \otimes \chi_2 \in \widehat{T}(L)$ 是经典权 $k$ 的, 则 $\bigl(\pind{}_{B(\Qp)}^{\GL_2(\Qp)} \chi_1 \otimes \chi_2\bigr)(\Qp)^{\mathrm{lp},\le k-2}$ 是 $\bigl(\pind{}_{B(\Qp)}^{\GL_2(\Qp)} \chi_1 \otimes \chi_2\bigr)(\Qp)^{\mathrm{lp}}$ 中被 $(\mathfrak{gl}_2(\Qp),B(\Qp))$ 稳定的子空间, \eqref{eq:2} 的像等于这个子空间, 而 \eqref{eq:2} 的核同构于:
\[
X_{-}^{k-1} \Bigl(U_{\mathfrak{gl}_2(\Qp)} \otimes_{U_{\mathfrak{b}(\Qp)}} \bigl(\pind{}_{B(\Qp)}^{\GL_2(\Qp)} \chi_1 \otimes \chi_2\bigr)(\Qp)^{\mathrm{lc}}\Bigr).
\]
\end{lemma}

\begin{proof}
保留前一引理证明的记号, 但这次假设 $w = k-2$ 是正整数或 $0$.
若 $L[z]_{\le w}$ 表示 $L[z]$ 中次数至多 $w$ 的多项式子空间, 由 \eqref{eq:3} 可见 $L[z]_{\le w}$ 在 $\mathfrak{gl}_2(\Qp)$ 作用下稳定.
更精确地, $L[z]_{\le w}$ 是 $\mathfrak{gl}_2(\Qp)$ 的不可约表示, 其最高权子空间 (关于 Borel 子代数 $\mathfrak{b}(\Qp)$) 恰是常值函数子空间 $L$, 在 $L$ 上 $\mathfrak{t}(\Qp)$ 通过主权 $\begin{pmatrix} \alpha & 0 \\ 0 & \delta \end{pmatrix} \mapsto w_2\alpha + w_1\delta$ 作用 (参见引理 \ref{lem:2.1.2} 的证明).
因此 $L$ 在 $U_{\mathfrak{gl}_2(\Qp)}$ 中的零化子等于 $X_{-}^{w+1}U_{\mathfrak{gl}_2(\Qp)}$, 而 \eqref{eq:4} 诱导一个满射:
\[
U_{\mathfrak{gl}_2(\Qp)} \otimes_{U_{\mathfrak{b}(\Qp)}} L \to L[z]_{\le w},
\]
其核为 $X_{-}^{w+1}U_{\mathfrak{gl}_2(\Qp)}$.
引理则从引理 \ref{lem:2.1.2} 证明中已用过的同一论证 (mutatis mutandis) 推出.
\end{proof}

下面的结果是"Amice--V\'elu 型"的, 是 \cite{ref:23} 的主要结果的一个特殊情形:

\begin{proposition}\label{prop:2.1.4}
若 $V$ 是 $L$ 上配备 $\GL_2(\Qp)$ 连续作用的 Banach 空间, 则同态的 (对偶) 限制诱导一个同构:
\[
\Hom_{\GL_2(\Qp)}\bigl(\anind_{B(\Qp)}^{\GL_2(\Qp)} \chi_1 \otimes \chi_2,\ V\bigr)
\xrightarrow{\sim}
\Hom_{(\mathfrak{gl}_2(\Qp),B(\Qp))}\bigl(\bigl(\pind{}_{B(\Qp)}^{\GL_2(\Qp)} \chi_1 \otimes \chi_2\bigr)(\Qp)^{\mathrm{lp}},\ V^{\mathrm{an}}\bigr),
\]
其中 $V^{\mathrm{an}}$ 是 $V$ 的局部解析向量子空间 (\cite[\S7]{ref:45}).
这里当然指连续同态 (注意, 由于 $\bigl(\pind{}_{B(\Qp)}^{\GL_2(\Qp)} \chi_1 \otimes \chi_2\bigr)(\Qp)^{\mathrm{lp}}$ 上的局部凸拓扑是最细的, 右边态射的连续性自动成立).
为方便读者, 我们在附录中给出这个重要而技术性的命题的完整证明.
\end{proposition}

最后回顾: 若 $\chi_1 \otimes \chi_2$ 是经典权 $k$ 的, 对表示 $\pind{\mathrm{lp},\le k-2}_{B(\Qp)}^{\GL_2(\Qp)} \chi_1 \otimes \chi_2$, 以及若 $\chi_1 \otimes \chi_2 \in \widehat{T}(L)$, 对 $\anind_{B(\Qp)}^{\GL_2(\Qp)} \chi_1 \otimes \chi_2$, 还可以联系它们的泛酉完备化, 即关于尽可能"细"的 $\GL_2(\Qp)$-不变连续半范数的完备化 (参见 \cite[\S1]{ref:21}).
它是一个酉 $\GL_2(\Qp)$-Banach 空间, 可能为零.
例如, 在局部代数情形, 它就是 $\pind{\mathrm{lp},\le k-2}_{B(\Qp)}^{\GL_2(\Qp)} \chi_1 \otimes \chi_2$ 关于 $\mathcal{O}_L[\GL_2(\Qp)]$ 上某个有限型 $\mathcal{O}_L$-格的完备化 (若这样的格存在).

\subsection{通过抛物线诱导的酉完备化定义 Banach 空间 $B(\sigma_p)$}\label{sec:2.2}

固定整数 $k \ge 2$ 和 $\chi_1 \otimes \chi_2 \in \widehat{T}(L)$, 使得 $\chi_1 \otimes \chi_2$ 是经典权 $k$ 的, $\chi_2$ 局部常值, 且 $\val(\chi_2(p)) = k-1$, $\val(\chi_1(p)) = 1-k$.
若 $k = 2$, 再假设 $\chi_1 \neq \chi_2|\,|^{2}$, 使得 $\lcind_{B(\Qp)}^{\GL_2(\Qp)} \chi_1 \otimes \chi_2$ 不可约且同构于 $\lcind_{B(\Qp)}^{\GL_2(\Qp)} \chi_2|\,| \otimes \chi_1|\,|^{-1}$.
设 $\eta_1 := \chi_1|\,|^{1-k}\varepsilon^{k-2} = \chi_1|\,|^{-1}z^{k-2}$ 和 $\eta_2 := \chi_2|\,|^{k-1}\varepsilon^{2-k} = \chi_2|\,|z^{2-k}$: 特征标 $\eta_1$ 和 $\eta_2$ 取整数值.

\begin{proposition}\label{prop:2.2.1}
局部代数表示 $\pind{\mathrm{lp},\le k-2}_{B(\Qp)}^{\GL_2(\Qp)} \chi_1 \otimes \chi_2$ 的泛酉完备化等同于可容许且拓扑不可约的酉 $\GL_2(\Qp)$-Banach 空间 $\cind_{B(\Qp)}^{\GL_2(\Qp)} \eta_2 \otimes \eta_1$.
\end{proposition}

\begin{proof}
利用缠结:
\[
\lcind_{B(\Qp)}^{\GL_2(\Qp)} \chi_1 z^{k-2} \otimes \chi_2 \simeq \lcind_{B(\Qp)}^{\GL_2(\Qp)} \chi_2|\,| \otimes \chi_1|\,|^{-1}z^{k-2},
\]
得到同构 $\pind{\mathrm{lp},\le k-2}_{B(\Qp)}^{\GL_2(\Qp)} \chi_1 \otimes \chi_2 \simeq \pind{\mathrm{lp},\le k-2}_{B(\Qp)}^{\GL_2(\Qp)} \eta_2 \otimes \eta_1$.
设 $I_{k-2}$ 为 $\Qp \to L$ 的函数 $f$ 的 $L$-向量空间, 它们在 $\Zp$ 之外为零, 且在 $\Zp$ 上的限制是次数 $\le k-2$ 的多项式, 并选择 $I_{k-2}$ 的取值于 $\mathcal{O}_L$ 的函数的 $\mathcal{O}_L$-格的一组基 $(f_0,\cdots,f_{k-2})$.
于是 $\bigoplus_{0 \le i \le k-2} \mathcal{O}_L[\GL_2(\Qp)]f_i \subset \pind{\mathrm{lp},\le k-2}_{B(\Qp)}^{\GL_2(\Qp)} \eta_2 \otimes \eta_1$ 这个子 $\GL_2(\Qp)$-表示等同于 $\pind{\mathrm{lp},\le k-2}_{B(\Qp)}^{\GL_2(\Qp)} \eta_2 \otimes \eta_1$ 中取值于 $\mathcal{O}_L$ 的函数的 $\GL_2(\Qp)$-稳定格.
因此可见此格是有限型的, 且 $\pind{\mathrm{lp},\le k-2}_{B(\Qp)}^{\GL_2(\Qp)} \eta_2 \otimes \eta_1$ (或 $\pind{\mathrm{lp},\le k-2}_{B(\Qp)}^{\GL_2(\Qp)} \chi_1 \otimes \chi_2$) 的泛酉完备化, 即 $\pind{\mathrm{lp},\le k-2}_{B(\Qp)}^{\GL_2(\Qp)} \eta_2 \otimes \eta_1$ 关于这个格的完备化, 等同于 $\cind_{B(\Qp)}^{\GL_2(\Qp)} \eta_2 \otimes \eta_1$.
后一个 $\GL_2(\Qp)$-Banach 空间可容许, 因为其对偶等同于 $L \otimes \bigl(\mathcal{O}_L[[\GL_2(\Zp)]] \otimes_{\mathcal{O}_L[[B(\Zp)]]} \mathcal{O}_L\bigr)$, 它在 $\mathcal{O}_L[[\GL_2(\Zp)]]$ 上有限型 (映射 $B(\Zp) \to \mathcal{O}_L$ 是 $\begin{pmatrix} a & b \\ 0 & d \end{pmatrix} \mapsto \eta_2(a)\eta_1(d)$).
不可约性可如下看出: 设 $B \subseteq \cind_{B(\Qp)}^{\GL_2(\Qp)} \eta_2 \otimes \eta_1$ 为被 $\GL_2(\Qp)$ 稳定的非零闭子空间.
由可容许 $p$ 进 Banach 空间的局部解析向量的稠密性 (\cite[Th.~7.1]{ref:45}), $B$ 包含 $\anind_{B(\Qp)}^{\GL_2(\Qp)} \eta_2 \otimes \eta_1$ 的一个在 $B$ 中稠密的子表示, 它本身必包含表示 $\pind{\mathrm{lp},\le k-2}_{B(\Qp)}^{\GL_2(\Qp)} \eta_2 \otimes \eta_1$ (因为后者局部多项式, 从而更是局部解析的).
但我们已经看到这个表示在 $\cind_{B(\Qp)}^{\GL_2(\Qp)} \eta_2 \otimes \eta_1$ 中稠密, 故后者也等于 $B$.
\end{proof}

回顾函数 $f : \Zp \to L$ 属于 $C^{k-1}$ 类, 若其 Mahler 展开 $f(z) = \sum_{n=0}^{+\infty} a_n(f)\binom{z}{n}$ 使得 $n^{k-1}|a_n(f)|$ 当 $n$ 趋于 $+\infty$ 时在 $\mathbb{R}_{+}$ 中趋于 $0$ (其中 $\binom{z}{0} = 1$, 若 $n > 0$ 则 $\binom{z}{n} = \frac{z(z-1)\cdots(z-n+1)}{n!}$).
记 $C^{k-1}(\Zp,L)$ 为这些函数的 $L$-向量空间.
对范数 $\|f\| := \mathrm{Sup}_n\, n^{k-1}|a_n(f)|$ 它是 Banach 空间.
也可以把此空间描述为 $k-1$ 次可微且 $(k-1)$ 次导数连续的函数 (\cite[\S54]{ref:42}).
$C^{k-1}$ 类函数对和、积、复合等运算是稳定的 (参见 \cite[Th.~77.5]{ref:42}).
设 $V$ 为函数 $f : \Qp \to L$ 的 $L$-向量空间, 使得函数 $f_1 := (z \mapsto f(pz))$ 和 $f_2 := (z \mapsto (\chi_2\chi_1^{-1})(z)f(1/z))$ 属于 $C^{k-1}(\Zp,L)$.
对范数 $\mathrm{Sup}(\|f_1\|,\|f_2\|)$ 它是 Banach 空间.
给 $V$ 配备 \eqref{eq:1} 给出的 $\GL_2(\Qp)$ 左作用, 这是连续自同构作用 (参见 \cite[\S4.2]{ref:3}).
对 $0 \le j \le k-2$ 和 $a \in \Qp$, 函数 $f(z) = z^{j}$ 和 $f(z) = (z-a)^{-j}(\chi_2\chi_1^{-1})(z-a)$ 属于 $V$ (证明类似于 \cite[Lem.~4.2.2]{ref:3}).
定义 $W \subset V$ 为 $V$ 中所有这些函数生成的子 $L$-向量空间在 $V$ 中的闭包.
它被 $\GL_2(\Qp)$ 稳定.

\begin{theorem}\label{thm:2.2.2}
局部解析表示 $\anind_{B(\Qp)}^{\GL_2(\Qp)} \chi_1 \otimes \chi_2$ 的泛酉完备化等同于商空间 $V/W$ 及其诱导的 $\GL_2(\Qp)$ 作用.
它是拓扑长度为 $2$ 的可容许酉 $\GL_2(\Qp)$-Banach 空间, 是 $\cind_{B(\Qp)}^{\GL_2(\Qp)} \eta_1\varepsilon \otimes \eta_2\varepsilon^{-1}$ 被 $\cind_{B(\Qp)}^{\GL_2(\Qp)} \eta_2 \otimes \eta_1$ 的非平凡扩张.
\end{theorem}

\begin{proof}
设 $\alpha := \chi_1(p)^{-1}$ 和 $\beta := p^{k-1}\chi_2(p)^{-1}$, 有:
\[
\anind_{B(\Qp)}^{\GL_2(\Qp)} \chi_1 \otimes \chi_2
= \varepsilon^{2-k} \otimes \anind_{B(\Qp)}^{\GL_2(\Qp)} \nr(\alpha^{-1})\psi_1 \otimes \nr(p\beta^{-1})\psi_2 z^{k-2},
\]
其中 $\psi_1 = \chi_1\nr(\chi_1(p))^{-1}\varepsilon^{k-2}$, $\psi_2 = \chi_2\nr(\chi_2(p))^{-1}$.
根据 \cite[Prop.~1.21]{ref:21} (更精确地, 是 loc.\ cit.\ 的证明), $\anind_{B(\Qp)}^{\GL_2(\Qp)} \chi_1 \otimes \chi_2$ 的泛酉完备化是关于由向量 $\mathbf{1}_{\Zp}(z)z^{j}$ 和 $\mathbf{1}_{\Qp \setminus \Zp}(z)(\chi_2\chi_1^{-1})(z)z^{-j}$ 生成的子 $\mathcal{O}_L[B(\Qp)]$-模的完备化, 其中 $\mathbf{1}_{U}$ 是开集 $U$ 的特征函数, $j \in \mathbb{Z}_{\ge 0}$.
通过与 \cite[Th.~4.3.1]{ref:3} 的证明中相同的计算 (此处我们把非分歧特征标 $\nr(\alpha^{-1})$ 和 $\nr(p\beta^{-1})$ 乘以整值特征标 $\psi_1$ 和 $\psi_2$, 并用整值特征标 $\varepsilon^{2-k}$ 扭曲表示, 但这不改变论证), 我们发现所寻求完备化的对偶 Banach 空间的一个 (以 $0$ 为中心的) 开球等同于 $\anind_{B(\Qp)}^{\GL_2(\Qp)} \chi_1 \otimes \chi_2$ 的强对偶中的分布 $\mu$, 使得对一切 $a \in \Qp$, $j \in \mathbb{Z}_{\ge 0}$ 和 $n \in \mathbb{Z}$:
\begin{align}
\label{eq:5}
&\int_{a+p^{n}\Zp} (z-a)^{j}\,d\mu(z) \in p^{n(j-k+1)}\mathcal{O}_L\\
\label{eq:6}
&\int_{\Qp - (a+p^{n}\Zp)} (z-a)^{-j}(\chi_2\chi_1^{-1})(z-a)\,d\mu(z) \in p^{n(k-1-j)}\mathcal{O}_L
\end{align}
(记住 $\val(\chi_1(p)) = 1-k$).
与 \cite[Th.~4.3.1]{ref:3} 证明类似的推理表明, 条件 \eqref{eq:5} 和 \eqref{eq:6} 恰好选出阶 $\le k-1$ 的缓增分布 (即 $V^{\vee}$ 的分布), 它们零化函数 $z^{j}$ 和 $(z-a)^{-j}(\chi_2\chi_1^{-1})(z-a)$ ($0 \le j \le k-2$) (在 \eqref{eq:5} 中令 $n \to -\infty$, 在 \eqref{eq:6} 中令 $n \to +\infty$), 即零化 $W$ 中的函数.
由此得到命题陈述的第一部分.
一个类似 (且更简单) 的推理表明, 诱导 $\anind_{B(\Qp)}^{\GL_2(\Qp)} \chi_1 z^{k-1} \otimes \chi_2 z^{1-k} = \anind_{B(\Qp)}^{\GL_2(\Qp)} \eta_1\varepsilon \otimes \eta_2\varepsilon^{-1}$ 的泛酉完备化等同于 $\cind_{B(\Qp)}^{\GL_2(\Qp)} \eta_1\varepsilon \otimes \eta_2\varepsilon^{-1}$, 且与命题 \ref{prop:2.2.1} 相同的证明表明后者是可容许且拓扑不可约的 $\GL_2(\Qp)$-Banach 空间.
由泛酉完备化的定义, 满射 $\anind_{B(\Qp)}^{\GL_2(\Qp)} \chi_1 \otimes \chi_2 \twoheadrightarrow \anind_{B(\Qp)}^{\GL_2(\Qp)} \chi_1 z^{k-1} \otimes \chi_2 z^{1-k}$ 诱导 $\GL_2(\Qp)$-等变满射 $V/W \twoheadrightarrow \cind_{B(\Qp)}^{\GL_2(\Qp)} \eta_1\varepsilon \otimes \eta_2\varepsilon^{-1}$.
由连续性, 这个满射仍等同于 $f \mapsto f^{(k-1)}$, 其核恰为 $N/W$, 其中 $N \subset V$ 是 $(k-1)$ 阶导数为零的函数的闭子空间 (注意 $W \subset N$).
另一方面, $\pind{\mathrm{lp},\le k-2}_{B(\Qp)}^{\GL_2(\Qp)} \chi_1 \otimes \chi_2$ 的泛酉完备化由函子性映入 $N/W$, 而通过命题 \ref{prop:2.2.1} 可见这必是闭浸入 (因为特征标 $\eta_1$ 和 $\eta_2$ 取整数值, 在 $\pind{\mathrm{lp},\le k-2}_{B(\Qp)}^{\GL_2(\Qp)} \chi_1 \otimes \chi_2 = \pind{\mathrm{lp},\le k-2}_{B(\Qp)}^{\GL_2(\Qp)} \eta_2 \otimes \eta_1$ 上相差等价只存在一个不变范数, 参见 e.g.\ \cite[Cor.~5.4.4]{ref:3}).
利用如下 $p$ 进分析结果: 次数至多 $k-2$ 的多项式函数子空间在 $C^{k-1}(\Zp,L)$ 中的闭包恰好等同于 $(k-1)$ 阶导数为零的函数 ($k = 2$ 情形的证明见 \cite[\S68]{ref:42}, 一般 $k$ 见 \cite[\S8]{ref:41}), 我们确实得到 $\GL_2(\Qp)$-等变正合列:
\[
0 \to \cind_{B(\Qp)}^{\GL_2(\Qp)} \eta_2 \otimes \eta_1 \to V/W \to \cind_{B(\Qp)}^{\GL_2(\Qp)} \eta_1\varepsilon \otimes \eta_2\varepsilon^{-1} \to 0,
\]
它 (利用前述结果和命题 \ref{prop:2.2.1}) 使 $V/W$ 成为拓扑长度为 $2$ 的可容许 $\GL_2(\Qp)$-Banach 空间.
最后, 此正合列非分裂, 否则局部解析表示 $\anind_{B(\Qp)}^{\GL_2(\Qp)} \chi_1 \otimes \chi_2$ 也会分裂, 而事实并非如此 (参见 \cite{ref:44}).
\end{proof}

\subsection{表示 $B(\sigma_p)$ 的一般构造与 $k = 1$, $k = 2$, $\eta_1 = \eta_2$ 的特殊情形}\label{sec:2.3}

设 $\sigma_p$ 是 $\Gal(\overline{\Qp}/\Qp)$ 的 $L$ 上二维可约潜在晶体表示.
可以写成:
\[
\sigma_p \simeq \begin{pmatrix} \chi_1|\,|^{1-k}\varepsilon^{k-2} & * \\ 0 & \chi_2|\,|^{k-1}\varepsilon^{1-k} \end{pmatrix} \otimes \eta
\]
其中 $\eta : \Gal(\overline{\Qp}/\Qp) \to L^{\times}$ 是连续特征标, $k \ge 1$ 是唯一确定的整数, $\chi_1 \otimes \chi_2$ 是经典权 $k$ 的特征标, $\chi_2$ 局部常值, 且 $\val(\chi_2(p)) = k-1$, $\val(\chi_1(p)) = 1-k$ (于是, 通过局部互反律, $\chi_1|\,|^{1-k}$ 和 $\chi_2|\,|^{k-1}$ 可延拓到 $\Gal(\overline{\Qp}/\Qp)$).
注意当 $k = 1$ 时,"经典权 $1$"这个术语是记号上的滥用, 这里的意思是 $\chi_2/\chi_1$ 是一个局部常值特征标与特征标 $z \mapsto z^{-1}$ 的乘积 (参见定义 \ref{def:2.1.1}).
若扩张 $*$ 非零, 则它唯一.
如同 \ref{sec:2.2} 设 $\eta_1 := \chi_1|\,|^{1-k}\varepsilon^{k-2}$ 和 $\eta_2 := \chi_2|\,|^{k-1}\varepsilon^{2-k}$.
若 $k = 2$, 假设 $\eta_1 \neq \eta_2$ (关于 $\eta_1 = \eta_2$ 的情形参见注记 \ref{rem:2.3.1}).
于是我们给 $\sigma_p$ "联系"如下的酉 $\GL_2(\Qp)$-Banach 空间 $B(\sigma_p)$:
\begin{enumerate}
\item[(i)] 若 $\sigma_p \simeq \begin{pmatrix} \eta_1 & 0 \\ 0 & \eta_2\varepsilon^{-1} \end{pmatrix} \otimes \eta$, 则
\[
B(\sigma_p) := \cind_{B(\Qp)}^{\GL_2(\Qp)} \eta_2 \otimes \eta_1 \otimes \eta \;\oplus\; \cind_{B(\Qp)}^{\GL_2(\Qp)} \eta_1\varepsilon \otimes \eta_2\varepsilon^{-1} \otimes \eta,
\]
\item[(ii)] 若 $\sigma_p \simeq \begin{pmatrix} \eta_1 & * \\ 0 & \eta_2\varepsilon^{-1} \end{pmatrix} \otimes \eta$ 且 $* \neq 0$, 并设 $k \ge 2$ ($k = 1$ 参见注记 \ref{rem:2.3.2}), 则 $B(\sigma_p)$ 是局部解析抛物线诱导 $\anind_{B(\Qp)}^{\GL_2(\Qp)} \chi_1 \otimes \chi_2$ 的泛酉完备化被 $\eta$ 扭曲后所得.
由定理 \ref{thm:2.2.2}, 它是非分裂扩张:
\[
0 \to \cind_{B(\Qp)}^{\GL_2(\Qp)} \eta_2 \otimes \eta_1 \otimes \eta \to B(\sigma_p) \to \cind_{B(\Qp)}^{\GL_2(\Qp)} \eta_1\varepsilon \otimes \eta_2\varepsilon^{-1} \otimes \eta \to 0.
\]
\end{enumerate}

\begin{remark}\label{rem:2.3.1}
这里讨论 $k = 2$ 且 $\eta_1 = \eta_2$ 的情形.
注意此时 $\lcind_{B(\Qp)}^{\GL_2(\Qp)} \chi_1 \otimes \chi_2 = \lcind_{B(\Qp)}^{\GL_2(\Qp)} |\,| \otimes |\,|^{-1} \otimes \eta_1$ 不再同构于 $\lcind_{B(\Qp)}^{\GL_2(\Qp)} \chi_2|\,| \otimes \chi_1|\,|^{-1} = \lcind_{B(\Qp)}^{\GL_2(\Qp)} 1 \otimes 1 \otimes \eta_1$.
设 $B(2,\infty)$ 为 $\lcind_{B(\Qp)}^{\GL_2(\Qp)} |\,| \otimes |\,|^{-1}$ 的泛酉完备化 (它是平凡表示被 $St$ 的非分裂扩张, 其中如 \cite{ref:8} 那样, 记 $St$ 为 Steinberg 表示的泛酉完备化).
此处我们被诱使把与 $\sigma_p$ 相关联的酉 $\GL_2(\Qp)$-Banach 空间 $B(\sigma_p)$ 定义为:
\begin{enumerate}
\item[(i)] 若 $\sigma_p$ 分裂:
\[
B(\sigma_p) := B(2,\infty) \otimes \eta_1\eta \;\oplus\; \cind_{B(\Qp)}^{\GL_2(\Qp)} \varepsilon \otimes \varepsilon^{-1} \otimes \eta_1\eta,
\]
\item[(ii)] 若 $\sigma_p$ 不分裂, $B(\sigma_p)$ 定义为 $\anind_{B(\Qp)}^{\GL_2(\Qp)} |\,| \otimes |\,|^{-1} \otimes \eta_1\eta$ 的泛酉完备化, 且有非分裂扩张:
\[
0 \to B(2,\infty) \otimes \eta_1\eta \to B(\sigma_p) \to \cind_{B(\Qp)}^{\GL_2(\Qp)} \varepsilon \otimes \varepsilon^{-1} \otimes \eta_1\eta \to 0.
\]
\end{enumerate}
当然, 这个"非一般"情形在尖点模形式的整体语境中从不出现, 因而不能直接"检验"它.
然而, 间接的整体考虑应该能够检验非分裂情形, 我们下面来解释.
假设 $\sigma_p$ 不分裂 (且 $k = 2$, $\eta_1 = \eta_2$) 并假设, 为简化, $\eta_1 = \eta_2$ 和 $\eta$ 都是平凡特征标 (必要时扭曲).
表示 $\sigma_p$ 可以看作当 $L \to \infty$ (即 $\val(L) \to -\infty$) 时 \cite[Ex.~1.3.5]{ref:8} 的半稳定可约表示 $V(2,L)$ 的"极限" (在某种合适的意义下) (例如可以在一维射影空间中考虑这个极限, 该空间是平凡特征标被 $\varepsilon^{-1}$ 的非平凡扩张的同构类).
对 $L$ 的每个 (有限) 值, 考虑 \cite{ref:8} 中定义的酉 $\GL_2(\Qp)$-Banach 空间 $B(2,L)$.
那么在 $B(2,L)$ 的构造中取 $\val(L) \to -\infty$ 的极限非常容易, 并可验证得到上面的 $\GL_2(\Qp)$-Banach 空间 $B(2,\infty)$ (在 \cite[\S4.5]{ref:8} 的讨论中, 只需把表示 $\sigma(2,L)$ 换成在 $\sigma(2,L)$ 的定义中用 $\log L$ 代替 $\val$ 得到的表示 $\sigma(2,\infty)$).
把 $\infty$ 换成 $L$ 的有限值, 我们如上面的 (ii) 期望与 $V(2,L)$ 对应的酉 $\GL_2(\Qp)$-Banach 空间 $B(V(2,L))$ 给出一个短正合列:
\[
0 \to B(2,L) \to B(V(2,L)) \to \cind_{B(\Qp)}^{\GL_2(\Qp)} \varepsilon \otimes \varepsilon^{-1} \to 0
\]
(尽管目前尚不清楚如何构造这样的正合列).
上面提到的间接整体检验就在于取一个水平 $N$、权 $2$ 的新模形式, 使得 $\sigma_p(f) \simeq V(2,L) \otimes \eta$, 并问 (引言中定义的) $\Pi_p(f)$ 是否确实给出一个短正合列:
\[
0 \to B(2,L) \otimes \eta \to \Pi_p(f) \to \cind_{B(\Qp)}^{\GL_2(\Qp)} \varepsilon \otimes \varepsilon^{-1} \otimes \eta \to 0.
\]
目前只已知存在 $\GL_2(\Qp)$-等变嵌入 $B(2,L) \otimes \eta \hookrightarrow \Pi_p(f)$ (\cite{ref:9}).
\end{remark}

\begin{remark}\label{rem:2.3.2}
当 $k = 1$ 且 $* \neq 0$ 时, 我们不知道如何 (甚至是否能够) 构造 $\cind_{B(\Qp)}^{\GL_2(\Qp)} \eta_2 \otimes \eta_1$ 被 $\cind_{B(\Qp)}^{\GL_2(\Qp)} \eta_1\varepsilon \otimes \eta_2\varepsilon^{-1}$ 的非分裂扩张.
\end{remark}

\begin{remark}\label{rem:2.3.3}
正如扩张在 Galois 一侧唯一 (当它非零时), 我们可以问: 在可容许酉 $\GL_2(\Qp)$-Banach 空间的 Abel 范畴中, 除 $\anind_{B(\Qp)}^{\GL_2(\Qp)} \chi_1 \otimes \chi_2$ 的泛酉完备化给出的那个之外, 是否还存在 $\cind_{B(\Qp)}^{\GL_2(\Qp)} \eta_2 \otimes \eta_1$ 被 $\cind_{B(\Qp)}^{\GL_2(\Qp)} \eta_1\varepsilon \otimes \eta_2\varepsilon^{-1}$ 的其它扩张.
\end{remark}

\begin{remark}\label{rem:2.3.4}
\ref{sec:2.2} 和 \ref{sec:2.3} 的结果与定义的一个版本已见于 \cite[\S7]{ref:4} (它不会发表).
\end{remark}
